% !TEX program = xelatex
% =====================================================================
%  SLIDE BẢO VỆ ĐỒ ÁN TRÍ TUỆ NHÂN TẠO, HỌC PHẦN AIP202
%  Đề tài: FAMS — Điểm danh bằng nhận diện khuôn mặt
%
%  Build:  latexmk -xelatex slide.tex
%  Overleaf: Menu > Compiler > XeLaTeX, Main document > slide.tex
%
%  Slide dùng chung số liệu với quyển báo cáo: cùng nạp auto/du-lieu-du-an.tex
%  do tools/quet-du-an.py sinh ra, nên hai tài liệu không bao giờ lệch nhau.
% =====================================================================
\documentclass[aspectratio=169, 11pt]{beamer}
\usepackage{styles/slide-dau}

% Số liệu quét từ kho mã nguồn. Sinh lại: python tools/quet-du-an.py
% File nay do baocao/tools/quet-du-an.py sinh ra. KHONG SUA TAY.
% Lan quet gan nhat: 12/11/2026 tren commit b099462.
% Sinh lai:  python baocao/tools/quet-du-an.py

\newcommand{\FsoEndpoint}{100}
\newcommand{\FsoModuleApi}{16}
\newcommand{\FsoBang}{21}
\newcommand{\FsoMigration}{15}
\newcommand{\FsoVaiTro}{9}
\newcommand{\FsoTest}{278}
\newcommand{\FsoFileTest}{27}
\newcommand{\FsoManHinh}{34}
\newcommand{\FsoNhomManHinh}{8}
\newcommand{\FsoDongPython}{6.649}
\newcommand{\FsoDongDart}{11.190}
\newcommand{\FsoDongTest}{4.096}
\newcommand{\FcommitCuoi}{b099462}
\newcommand{\FnhanhGit}{feature/flutter-phan-quyen}
\newcommand{\FsoCommit}{31}
\newcommand{\FngayCommit}{05/11/2026}
\newcommand{\FngayQuet}{12/11/2026}
\newcommand{\FchieuVector}{512}
\newcommand{\FnguongNhanDien}{0{,}45}
\newcommand{\FnguongTrung}{0{,}6}
\newcommand{\FnguongThat}{0{,}5}
\newcommand{\FsoKhungXacNhan}{2}
\newcommand{\FphutMuon}{15}
\newcommand{\FphutMoSom}{15}
\newcommand{\FphutAnHan}{30}
\newcommand{\FngayLuuAnh}{90}
\newcommand{\FphutHetHanUrl}{60}
\newcommand{\FphutToken}{480}
\newcommand{\FngayRefresh}{30}
\newcommand{\FsaiSoGps}{100}
\newcommand{\FphutOtp}{10}
\newcommand{\FsoLanSaiToiDa}{5}
\newcommand{\FbanKinhPhong}{50}
\newcommand{\FgiayDoiQr}{30}
\newcommand{\FsoLanGuiEmail}{3}
\newcommand{\FphienBanDongY}{2026-10}
\newcommand{\FtiLeCanhBao}{20}


% >>> SINH VIÊN ĐIỀN: họ tên, MSSV, lớp
\sdSinhVien{Nguyễn Đăng Sen}
\sdMSSV{2351220206}
\sdLop{CNTT 23A}                % <<< KIỂM TRA LẠI: đang đoán theo khoá 23
\sdGVHD{TS. Đỗ Phúc Hảo}
\sdMonHoc{Học phần AIP202 \textperiodcentered{} 2 tín chỉ}
\sdNoiNgay{Đà Nẵng, tháng \dots{} năm 2026}

\title{FAMS — Điểm danh bằng nhận diện khuôn mặt}
\author{\sdTenSinhVien}
\date{\sdDiaDiemNgay}

\begin{document}

\slidebia{HỆ THỐNG QUẢN LÝ SINH VIÊN TÍCH HỢP\\ĐIỂM DANH BẰNG NHẬN DIỆN KHUÔN MẶT}

% ---------------------------------------------------------------------
\begin{frame}{Nội dung trình bày}
  \begin{columns}[T]
    \begin{column}{0.5\textwidth}
      \begin{enumerate}
        \item Bài toán và mục tiêu
        \item Cơ sở lý thuyết
        \item Thiết kế hệ thống
      \end{enumerate}
    \end{column}
    \begin{column}{0.5\textwidth}
      \begin{enumerate}
        \setcounter{enumi}{3}
        \item Chống gian lận và bảo vệ dữ liệu
        \item Kết quả và demo
        \item Hạn chế, hướng phát triển
      \end{enumerate}
    \end{column}
  \end{columns}
  \vfill
  \begin{center}
    \solon{\FsoEndpoint}{endpoint} \hfill
    \solon{\FsoBang}{bảng dữ liệu} \hfill
    \solon{\FsoTest}{ca kiểm thử}
  \end{center}
\end{frame}

% =====================================================================
\section{Bài toán và mục tiêu}

\begin{frame}{Điểm danh thủ công: ba vấn đề}
  \begin{columns}[T]
    \begin{column}{0.33\textwidth}
      \begin{block}{Tốn thời gian}
        5--10 phút mỗi tiết để gọi tên hoặc chuyền tờ ký.
      \end{block}
    \end{column}
    \begin{column}{0.33\textwidth}
      \begin{block}{Không xác thực}
        Ký thay, dạ thay. Không để lại bằng chứng kiểm chứng được.
      \end{block}
    \end{column}
    \begin{column}{0.33\textwidth}
      \begin{block}{Tổng hợp quá muộn}
        Cuối kỳ mới biết ai vắng quá tỉ lệ, khi đã mất quyền thi.
      \end{block}
    \end{column}
  \end{columns}
  \vfill
  \begin{center}
    \small\color{gray700}
    Khuôn mặt luôn mang theo người, không cần phát thẻ, một khung hình
    nhận được nhiều người cùng lúc.
  \end{center}
\end{frame}

\begin{frame}{Năm mục tiêu của đồ án}
  \begin{enumerate}
    \item Quy trình nhận diện hoàn chỉnh, chạy được trên máy chủ \textbf{không GPU}.
    \item Chống ba kiểu gian lận: ảnh in, nhận nhầm, điểm danh từ xa.
    \item Số hoá trọn vẹn nghiệp vụ chuyên cần, có lưu vết mọi lần sửa.
    \item Phân quyền đúng cơ cấu tổ chức của trường, \FsoVaiTro{} vai trò.
    \item Một mã nguồn cho Web, Android, iOS, kèm bộ kiểm thử tự động.
  \end{enumerate}
  \vfill
  {\small\color{gray500}Toàn bộ mục tiêu được đối chiếu lại ở phần Kết quả.}
\end{frame}

% =====================================================================
\section{Cơ sở lý thuyết}

\begin{frame}{Bốn bước của quy trình nhận diện}
  \begin{center}
  \begin{tikzpicture}[scale=0.8, transform shape, node distance=0.7cm and 0.7cm, every node/.append style={align=center}]
    \node (anh)   [objbox] {Khung hình};
    \node (detect)[process, right=of anh]    {Phát hiện\\ SCRFD};
    \node (align) [process, right=of detect] {Căn chỉnh\\ 5 điểm mốc};
    \node (embed) [process, right=of align]  {ArcFace\\ \FchieuVector{} chiều};
    \node (match) [decision, right=of embed] {Cosine\\ $\geq \FnguongNhanDien$?};
    \node (out)   [startstop, right=of match] {Danh tính};
    \draw [flowarr] (anh) -- (detect);
    \draw [flowarr] (detect) -- (align);
    \draw [flowarr] (align) -- (embed);
    \draw [flowarr] (embed) -- (match);
    \draw [flowarr] (match) -- (out);
  \end{tikzpicture}
  \end{center}
  \vfill
  \begin{itemize}
    \item Không phân loại trực tiếp, mà học \textbf{hàm nhúng}: cùng người thì gần nhau.
    \item Thêm một sinh viên mới chỉ cần \textbf{một ảnh}, không huấn luyện lại.
  \end{itemize}
\end{frame}

\begin{frame}{ArcFace và ngưỡng nhận diện}
  \begin{columns}[T]
    \begin{column}{0.52\textwidth}
      Biên góc cộng tính kéo các mẫu của cùng một người lại gần tâm lớp
      chặt hơn mức tối thiểu cần để phân loại đúng:
      \vspace{0.3em}
      \[
        L = -\frac{1}{N}\sum_{i}
        \log \frac{e^{s\cos(\theta_{y_i}+m)}}
        {e^{s\cos(\theta_{y_i}+m)} + \sum_{j \neq y_i} e^{s\cos\theta_j}}
      \]
      \vspace{0.1em}
      Vector đã chuẩn hoá nên so khớp chỉ còn một phép nhân vô hướng.
    \end{column}
    \begin{column}{0.44\textwidth}
      \begin{block}{Hai loại lỗi không cân xứng}
        \textbf{Chấp nhận sai}: điểm danh cho người vắng mặt, phá hỏng cả hệ thống.

        \vspace{4pt}
        \textbf{Từ chối sai}: chỉ gây phiền, giảng viên xác nhận tay được.
      \end{block}
      \vspace{0.2em}
      {\small Ngưỡng $\tau = \FnguongNhanDien$ được chọn nghiêng về phía
      giảm chấp nhận sai.}
    \end{column}
  \end{columns}
\end{frame}

% =====================================================================
\section{Thiết kế hệ thống}

\begin{frame}{Kiến trúc ba lớp}
  \begin{center}
  \begin{tikzpicture}[scale=0.95, transform shape, node distance=0.55cm and 0.9cm, every node/.append style={align=center}]
    \node (web)  [objbox] {Flutter Web\\ cán bộ, giảng viên};
    \node (mobi) [objbox, below=of web] {Flutter Android, iOS\\ sinh viên, phụ huynh};
    \node (api)  [clsbox, right=of web, yshift=-0.8cm, minimum height=1.5cm]
          {FastAPI\\ \FsoEndpoint{} endpoint\\ xác thực, phân quyền};
    \node (svc)  [clsbox, right=of api, minimum height=1.5cm]
          {Lớp dịch vụ\\ nhận diện, chống giả mạo,\\ phiên, vị trí, thống kê};
    \node (db)   [membox, right=of svc, yshift=0.45cm] {PostgreSQL 16\\ pgvector, HNSW};
    \node (onnx) [membox, right=of svc, yshift=-0.55cm] {ONNX Runtime\\ SCRFD, ArcFace,\\ MiniFASNetV2};
    \draw [flowarr] (web) -- (api);
    \draw [flowarr] (mobi) -- (api);
    \draw [flowarr] (api) -- (svc);
    \draw [flowarr] (svc) -- (db);
    \draw [flowarr] (svc) -- (onnx);
  \end{tikzpicture}
  \end{center}
  \vfill
  {\small Mọi suy luận mô hình nằm ở máy chủ. Ba vòng lặp nền dọn ảnh quá
  hạn, tự đóng phiên bị quên và gửi email thông báo cho phụ huynh.}
\end{frame}

\begin{frame}{Phân quyền hai chiều}
  \begin{center}
  \footnotesize
  \setlength{\tabcolsep}{4pt}
  \begin{tabular}{@{}lcccccc@{}}
    \toprule
    Chức năng & Quản trị & P.\,ĐT & P.\,CTSV & LĐ Khoa & VP Khoa & Ban ngành \\
    \midrule
    Lớp học phần, lịch, ghi danh & Sửa   & Sửa & Xem   & Xem*   & Sửa* & -- \\
    Hồ sơ sinh viên, phụ huynh   & Sửa   & Xem & Sửa   & Xem*   & Xem* & -- \\
    Đồng ý bản giấy, duyệt mặt   & Sửa   & --  & Sửa   & --     & --   & -- \\
    Đặt lại mật khẩu SV, phụ huynh & Sửa & --  & Sửa   & --     & --   & -- \\
    Duyệt nghỉ có phép           & Duyệt & --  & Duyệt & Duyệt* & --   & -- \\
    Thống kê tổng hợp            & Xem   & Xem & Xem   & Xem*   & Xem* & Xem \\
    Nhật ký hệ thống             & Xem   & --  & --    & Xem*   & --   & -- \\
    \bottomrule
  \end{tabular}
  \end{center}
  \vfill
  \begin{itemize}
    \item \textbf{Quyền chức năng} kết hợp \textbf{phạm vi dữ liệu}: dấu sao
          nghĩa là chỉ trong Khoa của mình.
    \item Cài đặt tập trung thành bộ lọc truy vấn, mặc định là \textbf{từ chối}.
  \end{itemize}
\end{frame}

% =====================================================================
\section{Chống gian lận và bảo vệ dữ liệu}

\begin{frame}{Sáu lớp phòng vệ xếp chồng}
  \begin{center}
  \small
  \begin{tabularx}{0.95\textwidth}{@{}>{\bfseries\raggedright\arraybackslash}p{4.6cm} X@{}}
    \toprule
    Lớp phòng vệ & Ngăn kiểu gian lận \\
    \midrule
    Kiểm tra chất lượng ảnh & Ảnh mờ, tối, khuôn mặt quá nhỏ \\
    Chống giả mạo MiniFASNetV2 & Ảnh in, ảnh trên màn hình điện thoại \\
    Xác nhận \FsoKhungXacNhan{} khung hình & Một khung hình may mắn vượt ngưỡng \\
    So khớp theo phạm vi lớp & Nhận nhầm người ngoài lớp, chụp hộ bạn \\
    Đối chiếu GPS theo phòng học & Điểm danh từ nhà \\
    Mã QR đổi mỗi \FgiayDoiQr{} giây & Chụp mã gửi cho bạn ở xa \\
    \bottomrule
  \end{tabularx}
  \end{center}
  \vfill
  {\small\color{gray700} Không lớp nào là tuyệt đối. Một hành vi gian lận
  phải vượt qua \textbf{toàn bộ} mới thành công.}
\end{frame}

\begin{frame}{Dữ liệu khuôn mặt là dữ liệu sinh trắc}
  \begin{columns}[T]
    \begin{column}{0.33\textwidth}
      \begin{block}{Đồng ý là bắt buộc}
        Nội dung có phiên bản; tự xác nhận hoặc ký bản giấy. Chưa đồng ý thì
        không nhận diện.
      \end{block}
    \end{column}
    \begin{column}{0.33\textwidth}
      \begin{block}{Giữ tối thiểu}
        Chỉ lưu vector. Ảnh duyệt xoá ngay, ảnh bằng chứng xoá sau
        \FngayLuuAnh{} ngày.
      \end{block}
    \end{column}
    \begin{column}{0.33\textwidth}
      \begin{block}{Người dùng kiểm soát}
        Rút đồng ý là xoá ngay toàn bộ ảnh và vector khuôn mặt.
      \end{block}
    \end{column}
  \end{columns}
  \vfill
  \begin{itemize}
    \item Ảnh phát qua \textbf{đường dẫn được ký điện tử} (HMAC), hết hạn sau
          \FphutHetHanUrl{} phút, không có thư mục ảnh công khai.
    \item Không endpoint nào trả vector đặc trưng ra ngoài.
  \end{itemize}
\end{frame}

\begin{frame}{Nhật ký hệ thống: ai làm gì, lúc nào}
  \begin{columns}[T]
    \begin{column}{0.33\textwidth}
      \begin{block}{Ghi gì}
        Đăng nhập, cả lần sai. Mọi thao tác thêm, sửa, xoá. Xuất báo cáo.
      \end{block}
    \end{column}
    \begin{column}{0.33\textwidth}
      \begin{block}{Ghi thế nào}
        Tự động ở tầng ánh xạ đối tượng, trong cùng giao dịch: không
        endpoint nào quên ghi.
      \end{block}
    \end{column}
    \begin{column}{0.33\textwidth}
      \begin{block}{Ai xem}
        Quản trị viên toàn trường; lãnh đạo Khoa phần của Khoa mình.
      \end{block}
    \end{column}
  \end{columns}
  \vfill
  \begin{itemize}
    \item Mỗi dòng lưu giá trị cũ và mới của từng trường đã đổi.
    \item Mật khẩu, vector khuôn mặt, mã OTP chỉ ghi là \emph{đã thay đổi}.
  \end{itemize}
\end{frame}

% =====================================================================
\section{Kết quả}

\begin{frame}{Sản phẩm thu được}
  \begin{center}
    \solon{\FsoEndpoint}{endpoint} \hfill
    \solon{\FsoBang}{bảng dữ liệu} \hfill
    \solon{\FsoManHinh}{màn hình ứng dụng}

    \vspace{0.7cm}

    \solon{\FsoVaiTro}{vai trò người dùng} \hfill
    \solon{\FsoTest}{ca kiểm thử} \hfill
    \solon{\FsoCommit}{commit}
  \end{center}
  \vfill
  {\small\color{gray500} Mọi con số trên được quét tự động từ kho mã nguồn,
  bản commit \code{\FcommitCuoi} ngày \FngayCommit.}

  \vspace{4pt}
  {\small\bfseries Mã nguồn:
  \href{https://github.com/jfgafgfagfbja/FAMS}{\color{daumaroon}github.com/jfgafgfagfbja/FAMS}}
\end{frame}

\begin{frame}{Sinh viên tự điểm danh: mã QR, vị trí, khuôn mặt}
  \begin{columns}[c]
    \begin{column}{0.34\textwidth}
      \centering
      \includegraphics[width=\linewidth]{kq-ma-qr.png}
    \end{column}
    \begin{column}{0.64\textwidth}
      \centering
      \includegraphics[width=0.325\linewidth]{kq-tu-dd-1.png}\hfill
      \includegraphics[width=0.325\linewidth]{kq-tu-dd-2.png}\hfill
      \includegraphics[width=0.325\linewidth]{kq-tu-dd-3.png}
    \end{column}
  \end{columns}
  \vfill
  {\small\color{gray500} Chạy thật trên iPhone: quét mã, đối chiếu vị trí, rồi so khớp chỉ với khuôn mặt của chính sinh viên.}
\end{frame}

\begin{frame}{Sinh viên và phụ huynh trên điện thoại}
  \begin{center}
    \includegraphics[height=0.72\textheight]{kq-sv-trang-chu.png}\hfill
    \includegraphics[height=0.72\textheight]{kq-sv-chuyen-can.png}\hfill
    \includegraphics[height=0.72\textheight]{kq-sv-khuon-mat.png}\hfill
    \includegraphics[height=0.72\textheight]{kq-ph-thong-bao.png}
  \end{center}
  \vfill
  {\small\color{gray500} Trang chủ, chuyên cần, dữ liệu khuôn mặt của sinh viên; thông báo của phụ huynh.}
\end{frame}

\begin{frame}{Phòng ban: tổng quan chuyên cần}
  \begin{columns}[c]
    \begin{column}{0.49\textwidth}
      \includegraphics[width=\linewidth]{kq-tong-quan.png}
    \end{column}
    \begin{column}{0.49\textwidth}
      \includegraphics[width=\linewidth]{kq-bieu-do.png}
    \end{column}
  \end{columns}
  \vfill
  {\small\color{gray500} Thông tin cán bộ, lối tắt theo vai trò và biểu đồ theo Khoa. Số liệu từ dữ liệu mẫu.}
\end{frame}

\begin{frame}{Kiểm thử tự động}
  \begin{columns}[T]
    \begin{column}{0.5\textwidth}
      \begin{itemize}
        \item \FsoTest{} ca trong \FsoFileTest{} tệp, khoảng \FsoDongTest{}
              dòng mã kiểm thử.
        \item Chạy trên \textbf{PostgreSQL thật}: kiểu vector không mô phỏng
              được bằng cơ sở dữ liệu nhúng.
        \item Mỗi ca chạy trong một giao dịch rồi quay lui.
      \end{itemize}
    \end{column}
    \begin{column}{0.46\textwidth}
      \begin{block}{Nhóm kiểm thử trọng tâm}
        \small
        Phân quyền theo vai trò và theo Khoa \\
        Vòng đời phiên điểm danh \\
        Duyệt khuôn mặt và quyền riêng tư \\
        Tự điểm danh kèm đối chiếu GPS \\
        Token, thu hồi và đường dẫn ảnh ký điện tử \\
        Đồng ý sinh trắc và nhật ký hệ thống
      \end{block}
    \end{column}
  \end{columns}
  \vfill
  {\small\color{gray700} Lần chạy ngày 12/11/2026: \code{278 passed in
  39.56s}.}
\end{frame}

\begin{frame}{Demo}
  \begin{center}
    \vspace{0.6cm}
    {\Large\bfseries\color{daumaroon} Phần trình diễn trực tiếp}\\[0.7cm]
    \begin{tabular}{@{}r@{\hspace{0.5cm}}l@{}}
      \bfseries 1 & Sinh viên đăng ký khuôn mặt, cán bộ duyệt \\[3pt]
      \bfseries 2 & Giảng viên mở phiên và quét camera cả lớp \\[3pt]
      \bfseries 3 & Thử gian lận bằng ảnh in, hệ thống chặn \\[3pt]
      \bfseries 4 & Sinh viên tự điểm danh, đối chiếu vị trí \\[3pt]
      \bfseries 5 & Cảnh báo vắng và xuất báo cáo Excel \\
    \end{tabular}
    \vspace{0.5cm}

    {\small\color{gray500} Tài khoản mẫu cho từng vai trò có trong phụ lục
    của quyển báo cáo.}
  \end{center}
\end{frame}

% ---------------------------------------------------------------------
% ẢNH DEMO: chụp màn hình hệ thống, đặt vào figs/ rồi bỏ ký tự % ở đầu.
% ---------------------------------------------------------------------
% \begin{frame}{Màn hình quét camera của giảng viên}
%   \begin{center}
%     \includegraphics[height=0.72\textheight]{man-quet-camera.png}
%   \end{center}
% \end{frame}

% =====================================================================
\section{Hạn chế và hướng phát triển}

\begin{frame}{Điều quan trọng nhất học được}
  \vfill
  \begin{center}
    {\Large\bfseries\color{daumaroon}
     Phần khó của một hệ thống trí tuệ nhân tạo\\ nằm ở \emph{ngoài} mô hình.}
  \end{center}
  \vfill
  \begin{itemize}
    \item Ghép SCRFD và ArcFace vào hệ thống là phần tốn ít công nhất.
    \item Phần tốn công là những câu hỏi mô hình không trả lời được: phiên
          đóng lúc nào, ai được sửa điểm danh, sửa thì lưu vết ra sao, sinh
          viên rút lại dữ liệu thì xoá những gì.
    \item Ràng buộc nên đặt ở tầng thấp nhất có thể, tức ở cơ sở dữ liệu.
  \end{itemize}
\end{frame}

\begin{frame}{Hạn chế và hướng phát triển}
  \begin{columns}[T]
    \begin{column}{0.48\textwidth}
      \textbf{\color{daumaroon} Hạn chế}
      \begin{itemize}
        \item Chưa kiểm định trên tập dữ liệu lớn.
        \item Ngưỡng nhận diện cố định cho mọi người.
        \item Chống giả mạo chỉ dựa trên ảnh tĩnh.
        \item Bộ đếm xác nhận mất khi khởi động lại máy chủ.
        \item Chưa chạy thử ở quy mô thật.
      \end{itemize}
    \end{column}
    \begin{column}{0.48\textwidth}
      \textbf{\color{agentragreen} Hướng phát triển}
      \begin{itemize}
        \item Đánh giá trên tập lớn, chọn ngưỡng theo đường cong ROC.
        \item Ngưỡng thích nghi theo từng sinh viên.
        \item Bổ sung tín hiệu chiều sâu khi thiết bị hỗ trợ.
        \item Đưa bộ đếm ra kho lưu trữ ngoài tiến trình.
        \item Kết nối hệ thống quản lý đào tạo, thông báo đẩy.
      \end{itemize}
    \end{column}
  \end{columns}
\end{frame}

\begin{frame}[plain]
  \vfill
  \begin{center}
    {\fontsize{26}{30}\selectfont\bfseries\color{daumaroon} Cảm ơn thầy cô đã lắng nghe}\\[6pt]
    {\color{daugold}\rule{6cm}{1.6pt}}\\[0.8cm]
    {\large\color{daunavy} Em xin nhận câu hỏi của hội đồng}\\[1.0cm]
    {\small\color{gray500}
     \sdTenSinhVien{} \textperiodcentered{} \sdMaSinhVien{} \textperiodcentered{}
     \sdTenLop \\[3pt]
     Giảng viên hướng dẫn: \sdTenGVHD}
  \end{center}
  \vfill
\end{frame}

\end{document}
